\documentclass[tikz,border=10pt]{standalone}
\usepackage[no-math]{fontspec}
\usepackage[brazil]{babel}
\usepackage{amsmath,amssymb}
\IfFileExists{C:/Windows/Fonts/Lato_Regular.ttf}{%
  \setmainfont{Lato}[
    Path = C:/Windows/Fonts/,
    UprightFont = Lato_Regular.ttf,
    BoldFont = Lato_Bold.ttf,
    ItalicFont = Lato_Italic.ttf,
    BoldItalicFont = Lato_BoldItalic.ttf
  ]
  \setsansfont{Lato}[
    Path = C:/Windows/Fonts/,
    UprightFont = Lato_Regular.ttf,
    BoldFont = Lato_Bold.ttf,
    ItalicFont = Lato_Italic.ttf,
    BoldItalicFont = Lato_BoldItalic.ttf
  ]
}{%
  \RequirePackage[default]{lato}
}
\renewcommand{\familydefault}{\sfdefault}

\usepackage{xcolor}
\usepackage{tikz}
\usepackage{pgfplots}
\pgfplotsset{compat=1.18}
\usetikzlibrary{arrows.meta,positioning,shapes.geometric,fit,backgrounds,calc,decorations.pathreplacing,patterns}

% Paleta institucional PLANTERR / UEFS (alinhada ao estilo planterr-slides.css / 09a)
\definecolor{etNavy}{HTML}{1F4268}
\definecolor{etNavyDark}{HTML}{152E4A}
\definecolor{etGold}{HTML}{B38A3D}
\definecolor{etGoldDark}{HTML}{8C6927}
\definecolor{etTeal}{HTML}{0E7C7B}
\definecolor{etForest}{HTML}{2D6A4F}
\definecolor{etAlert}{HTML}{9E2A2B}
\definecolor{etInk}{HTML}{26313D}
\definecolor{etMuted}{HTML}{5B6573}
\definecolor{etBorder}{HTML}{D5DEE8}

\definecolor{bgNavy}{HTML}{EDF2F7}
\definecolor{bgGold}{HTML}{F8F3E9}
\definecolor{bgTeal}{HTML}{EBF6F5}
\definecolor{bgForest}{HTML}{EEF5F0}
\definecolor{bgAlert}{HTML}{FDF2F2}
\definecolor{bgCanvas}{HTML}{FFFFFF}

\tikzset{
  >=Stealth,
  etBox/.style={
    draw=#1, fill=white, rounded corners=5pt, line width=1.15pt,
    inner sep=7pt, align=center, font=\footnotesize\sffamily, text=etInk
  },
  etTint/.style={
    draw=#1, fill=#1!8!white, rounded corners=5pt, line width=1.15pt,
    inner sep=7pt, align=center, font=\footnotesize\sffamily, text=etInk
  },
  etBadge/.style={
    fill=#1, text=white, rounded corners=9pt, inner xsep=8pt, inner ysep=3pt,
    font=\scriptsize\bfseries\sffamily
  },
  etTitle/.style={font=\normalsize\bfseries\sffamily, text=etNavy, align=center},
  etSub/.style={font=\footnotesize\sffamily, text=etMuted, align=center},
  etArrow/.style={->, >=Stealth, line width=1.35pt},
  etDArrow/.style={<->, >=Stealth, line width=1.35pt}
}

\usepgfplotslibrary{fillbetween}

\begin{document}
\begin{tikzpicture}[font=\sffamily]

\begin{axis}[
  width=15.0cm,
  height=8.4cm,
  title={\textbf{O ``Vale da Morte'' Tecnológico e a Transição Bancada $\rightarrow$ Mercado}},
  title style={font=\normalsize\bfseries\sffamily, text=etNavy, yshift=6pt},
  ylabel={Caixa acumulado (escala ilustrativa)},
  ylabel style={font=\footnotesize\sffamily, text=etInk},
  xmin=0, xmax=4.2,
  ymin=-4.5, ymax=5.8,
  xtick={0.55, 1.65, 2.85, 3.80},
  xticklabels={
    {Pesquisa /\\Ideação},
    {Protótipo /\\Validação},
    {Tração /\\Break-even},
    {Escala}
  },
  xticklabel style={align=center, font=\footnotesize\sffamily, text=etInk, yshift=-3pt},
  ytick={-4,-2,0,2,4},
  yticklabel style={font=\scriptsize\sffamily, text=etMuted},
  axis lines=left,
  axis line style={draw=etMuted, line width=0.9pt},
  tick style={draw=etMuted},
  ymajorgrids=true,
  grid style={draw=etBorder, dotted, line width=0.6pt},
  clip=false,
  legend pos=north west,
  legend style={
    draw=etBorder, fill=white, rounded corners=4pt,
    font=\scriptsize\sffamily, text=etInk,
    inner sep=5pt, cells={anchor=west}
  }
]

% Subtle vertical stage dividers
\draw[draw=etBorder, dashed, line width=0.7pt] (axis cs:1.10, -4.5) -- (axis cs:1.10, 5.5);
\draw[draw=etBorder, dashed, line width=0.7pt] (axis cs:2.25, -4.5) -- (axis cs:2.25, 5.5);
\draw[draw=etBorder, dashed, line width=0.7pt] (axis cs:3.40, -4.5) -- (axis cs:3.40, 5.5);

% Zero reference line
\addplot[name path=zero, domain=0:4.1, draw=etMuted, dashed, line width=0.95pt, forget plot] {0};

% Smooth cashflow curve: exact zero at x=0 and x=2.85, trough at x=1.65 (y=-3.2)
\addplot[
  name path=curve,
  smooth,
  tension=0.62,
  line width=2.2pt,
  etNavy
] coordinates {
  (0.00,  0.00)
  (0.55, -1.05)
  (1.15, -2.65)
  (1.65, -3.20)
  (2.15, -2.35)
  (2.85,  0.00)
  (3.45,  2.45)
  (4.05,  5.05)
};
\addlegendentry{Caixa acumulado}

% Shaded regions: Vale da Morte (< 0) and positive scale (> 0)
\addplot[area legend, fill=bgAlert, draw=etAlert!45!white, fill opacity=0.85]
  fill between[of=curve and zero, soft clip={domain=0:2.85}];
\addlegendentry{Vale da Morte (déficit acumulado)}

\addplot[area legend, fill=bgTeal, draw=etTeal!45!white, fill opacity=0.85]
  fill between[of=curve and zero, soft clip={domain=2.85:4.05}];
\addlegendentry{Retorno positivo e escala}

% Re-draw curve on top of shading so stroke stays sharp
\addplot[
  smooth,
  tension=0.62,
  line width=2.2pt,
  etNavy,
  forget plot
] coordinates {
  (0.00,  0.00)
  (0.55, -1.05)
  (1.15, -2.65)
  (1.65, -3.20)
  (2.15, -2.35)
  (2.85,  0.00)
  (3.45,  2.45)
  (4.05,  5.05)
};

% Key points
\fill[etAlert] (axis cs:1.65, -3.20) circle (3.3pt);
\fill[etTeal]  (axis cs:2.85,  0.00) circle (3.3pt);

% Non-overlapping Callout 1: Ponto de Mínimo (placed cleanly in open middle space above zero)
\node[
  draw=etAlert, fill=white, rounded corners=4pt, line width=1pt,
  inner sep=5pt, align=center, font=\scriptsize\sffamily, text=etInk
] (minBox) at (axis cs:1.65, 1.45) {%
  \textbf{\color{etAlert}Ponto de Mínimo}\\
  {\color{etMuted}(Maior risco financeiro)}
};
\draw[etArrow, etAlert, line width=1.15pt] (minBox.south) -- (axis cs:1.65, -3.02);

% Non-overlapping Callout 2: Break-even (placed cleanly to upper-left of zero crossing)
\node[
  draw=etTeal, fill=white, rounded corners=4pt, line width=1pt,
  inner sep=5pt, align=center, font=\scriptsize\sffamily, text=etInk
] (beBox) at (axis cs:2.65, 2.75) {%
  \textbf{\color{etTeal}Break-even (Payback)}\\
  {\color{etMuted}Transição ao equilíbrio}
};
\draw[etArrow, etTeal, line width=1.15pt] (beBox.south) -- (axis cs:2.83, 0.18);

\end{axis}
\end{tikzpicture}
\end{document}
